\section{Wie binde ich Literatur ein?}
\label{sec:literatur}

Literatur wird über eine \textit{literature.bib} Datei eingebunden. Aus dieser Datei können dann Quelle zitiert werden.

Diese Vorlage zitiert standardmäßig im \gls{ieee} style.

\begin{lstlisting}[caption=Beispiel Eintrag in literature.bib]
@article{einstein1908relativity,
  author = {Albert Einstein},
  journal = {Jahrbuch der Radioaktivität und Elektronik},
  title = {Über das Relativitätsprinzip und die aus demselben gezogenen Folgerungen},
  volume = {4},
  year = {1908},
  pages = {411--462},
}
\end{lstlisting}

Daraus kann dann mit \lstinline|\cite{einstein1908relativity}| zitiert werden:

"Es zeigte sich aber überraschenderweise, daß es nur nötig war, den Begriff der Zeit genügend scharf zu fassen, \dots" \cite{einstein1908relativity}

Ich Persönlich kann noch \href{https://www.zotero.org/}{Zotero} empfehlen um die Literatur Recherche zu machen, da dies auch einen export zu \textit{.bib} Dateien hat.